\documentclass[10pt]{article}
\usepackage[T1]{fontenc}
\usepackage[utf8]{inputenc}
\usepackage[french]{babel}
\usepackage{amsmath,amssymb,mathrsfs,mathpazo}
\usepackage{geometry}
\geometry{a4paper,margin=2cm}
\usepackage{graphicx,array,tabularx,tabularray,colortbl,multirow,multicol,booktabs,enumitem,pifont,ulem,fancybox,diagbox,cancel}
\usepackage{pstricks,pstricks-add,pst-plot,pst-tree,pst-node,pst-text,pst-func,pst-math,pst-bspline,pst-3dplot,pst-eucl,pst-solides3d}
\let\pscalculate\undefined
\usepackage{pst-bezier}
\usepackage{tikz,pgfplots,tkz-tab}
\usetikzlibrary{arrows,arrows.meta,patterns,calc,positioning,trees,decorations.markings,shapes.geometric,shapes.misc,shapes.arrows,intersections,decorations.pathmorphing,decorations.pathreplacing,backgrounds,shadows,babel}
\pgfplotsset{compat=1.18}
\usepgfplotslibrary{fillbetween}
\usepackage[np]{numprint}
\usepackage{listings,makecell,pas-tableur,twoopt}
\let\barre\undefined
\usepackage[pstricks]{tablvar}
\definecolor{darkgreen}{RGB}{0,100,0}
\providecommand{\red}{\color{red}}
\providecommand{\blue}{\color{blue}}
\providecommand{\green}{\color{green}}
\providecommand{\black}{\color{black}}
\providecommand{\white}{\color{white}}
\providecommand{\gray}{\color{gray}}
\providecommand{\cyan}{\color{cyan}}
\providecommand{\magenta}{\color{magenta}}
\providecommand{\yellow}{\color{yellow}}
\providecommand{\violet}{\color{violet}}
\providecommand{\orange}{\color{orange}}
\providecommand{\euro}{\texteuro}
\providecommand{\EUR}{\texteuro}
\providecommand{\ieme}{\textsuperscript{e}}
\providecommand{\degres}{\ensuremath{^\circ}}
\providecommand{\degree}{\ensuremath{^\circ}}
\providecommand{\vv}[1]{\overrightarrow{#1}}
\providecommand{\Coord}[1]{#1}
\providecommand{\decosix}{}
\providecommand{\starredbullet}{\ensuremath{\bullet}}
\pagestyle{empty}
\setlength{\parindent}{0pt}
\usepackage{adjustbox}\DeclareUnicodeCharacter{2236}{:}\DeclareUnicodeCharacter{2190}{\ensuremath{\leftarrow}}
\providecommand{\R}{\mathbb{R}}
\providecommand{\N}{\mathbb{N}}
\providecommand{\Z}{\mathbb{Z}}
\providecommand{\Q}{\mathbb{Q}}
\providecommand{\D}{\mathbb{D}}
\providecommand{\C}{\mathbb{C}}
\providecommand{\e}{{\mathrm{e}}}
\providecommand{\vect}[1]{\overrightarrow{#1}}
\providecommand{\vectt}[1]{\overrightarrow{\mathrm{#1}}}
\providecommand{\vvt}[1]{\overrightarrow{\mathrm{#1}}}
\providecommand{\np}[2][]{{#2~#1}}
\providecommand{\vv}[1]{{\overrightarrow{#1}}}
\renewcommand{\d}{{\mathrm{d}}}
\providecommand{\degres}{{^\circ}}
\providecommand{\degre}{{^\circ}}
\providecommand{\blue}{\color{blue}}
\providecommand{\red}{\color{red}}
\providecommand{\enskip}{\quad}
\providecommand{\llbracket}{\lbrack\!\lbrack}
\providecommand{\rrbracket}{\rbrack\!\rbrack}
\providecommand{\euro}{{\text{€}}}
\providecommand{\no}{{\text{n°}}}
\providecommand{\interval}[2]{{[#1;#2]}}
\providecommand{\intervFO}[2]{{[#1;#2[}}
\providecommand{\dotfill}{\ldots}

\providecommand{\ds}{\displaystyle}
\providecommand{\textding}[1]{#1}
\providecommand{\Oij}{\(\mathrm{O};\vec{\imath},\vec{\jmath}\)}
\providecommand{\Oijk}{\(\mathrm{O};\vec{\imath},\vec{\jmath},\vec{k}\)}
\providecommand{\Ouv}{\(\mathrm{O};\vec{u},\vec{v}\)}
\providecommand{\intervFF}[2]{\ensuremath{[#1;#2]}}
\renewcommand{\euro}{\ifmmode\text{\texteuro}\else\texteuro\fi}
\DeclareUnicodeCharacter{00B0}{\ifmmode^{\circ}\else\textdegree\fi}
\DeclareUnicodeCharacter{20AC}{\ifmmode\text{\texteuro}\else\texteuro\fi}
% Énoncés : Éducation nationale. Transcription et corrigés : APMEP, Denis Vergès et contributeurs. Découpage et classement : CoopMaths, Sébastien Lozano, Éric Elter et contributeurs. Adaptation de la présentation : Maths Explorer.
% CC BY-SA : https://coopmaths.fr/www/about/
% https://coopmaths.fr/alea/static/eam/2026/tex/eam_specifique_2026_06_metropole_automatismes.tex
\begin{document}
\noindent\textit{APMEP / CoopMaths. CC BY-SA. Présentation adaptée par Maths Explorer.}\par\medskip
\textbf{Pour cette première partie, aucune justification n'est demandée et une seule réponse est possible par question.}

	\textbf{Pour chaque question, reporter son numéro sur la copie et indiquer la réponse.}

	Une réponse fausse ou l'absence de réponse n'enlève aucun point.

\subsection*{Question 1}
	Le nombre \quad $\dfrac{2}{5}$ \quad est égal à :

	\begin{tblr}{@{}XXXX}
		\textbf{A.} 0,2&
		\textbf{B.} 0,25&
		\textbf{C.} 0,4&
		\textbf{D.} 0,5
	\end{tblr}

	\subsection*{Question 2}
	$30\,\%$ de 150 est égal à :

	\begin{tblr}{@{}XXXX}
		\textbf{A.} 15&
		\textbf{B.} 30&
		\textbf{C.} 45&
		\textbf{D.} 60
	\end{tblr}

	\subsection*{Question 3}
	\begin{minipage}{0.4\linewidth}
		On donne ci-contre la

		représentation graphique

		d'une fonction $f$.

	Lequel de ces nombres est un antécédent de 3 ?

	\begin{tblr}{@{}X[3]X[2]}
		\textbf{A.} 0,5&
		\textbf{B.} 1\\
		\textbf{C.} 1,5&
		\textbf{D.} 2
	\end{tblr}
	\end{minipage}
	\hfill
	\begin{tikzpicture}[baseline={(current bounding box.center)}]
		\draw [color = gray!75, line width = 0.5pt, xstep=0.5, ystep = 0.5] (-4.4,-2.4) grid (5.4,3.9);
		\draw (-4.4,0) -- (5.4,0);
		\foreach \x in {-4,...,-1,1,2,...,5}
		\draw[shift={(\x,0)},color=black] (0pt,2pt) -- (0pt,-2pt) node[below] {\footnotesize $\np{\x}$};
		\draw (0,-2.4) -- (0,3.9);
		\foreach \y in {-2,-1,1,2,3}
		\draw[shift={(0,\y)},color=black] (2pt,0pt) -- (-2pt,0pt) node[left] {\footnotesize $\np{\y}$};
		\draw[color=black] (-2pt,-2pt) node[below left] {\footnotesize 0};
		\draw[line width=1.2pt,red]
		(-4,-0.5) .. controls (-3.75,0.75) and (-3.3,1.5) ..
		(-3,1.5)  .. controls (-2.1 ,1.5 ) and (-2.4,-1.5) ..
		(-1.5,-1.5) .. controls (-1.25,-1.5 ) and (-1.1,-1.1) ..(-1,-1) --
		(0.5,2)  .. controls (0.6,2.15 ) and (0.8,3) ..
		(1,3)    .. controls (1.2,3) and (1.4,2.15 ) ..
		(1.5,2)  .. controls (1.6,1.8) and (2.4,1.2) ..
		(2.5,1)  .. controls (2.7,.8) and (2.8,0.5) ..
		(3,0.5)  .. controls (3.2,0.5) and (3.3,0.8) ..
		(3.5,1) -- (5,2.5);
	\end{tikzpicture}

	\subsection*{Question 4}
	La solution de l'équation \quad $7 x+4=5 x+6$ \quad est :

	\begin{tblr}{@{}XXXX}
		\textbf{A.} $x=-1$&
		\textbf{B.} $x=0$&
		\textbf{C.} $x=2$&
		\textbf{D.} $x=1$
	\end{tblr}

	\subsection*{Question 5}
	Un article coûte initialement $50$\,\euro{}.

	Son prix diminue de $10\,\%$ puis augmente de $10\,\%$.

	Son prix final est de :

	\begin{tblr}{@{}XXXX}
		\textbf{A.} 49,50 \euro{}&
		\textbf{B.} 49,90 \euro{}&
		\textbf{C.} 50 \euro{}&
		\textbf{D.} 50,10 \euro{}
	\end{tblr}

	\subsection*{Question 6}
	On considère la courbe d'équation \quad $y=2 x^{2}-x+3$.

	Le point d'abscisse $-1$ de la courbe est :

	\begin{tblr}{@{}XXXX}
		\textbf{A.} $\mathrm{A}(-1 ; 0)$&
		\textbf{B.} $\mathrm{B}(-1 ; 2)$&
		\textbf{C.} $\mathrm{C}(-1 ; 4)$&
		\textbf{D.} $\mathrm{D}(-1 ; 6)$
	\end{tblr}

	\subsection*{Question 7}
	Une droite passe par les points A$(-1 ; 2)$ et $\mathrm{B}(-3 ; 4)$.

	Son coefficient directeur est :

	\begin{tblr}{@{}XXXX}
		\textbf{A.} $-2$&
		\textbf{B.} $-1$&
		\textbf{C.} 1&
		\textbf{D.} 2
	\end{tblr}

	\subsection*{Question 8}
	On s'intéresse au confort d'un hôtel.

	Les six dernières notes obtenues sont : \quad $2 ; 3 ; 5 ; 4 ; 2 ; 3$

	La médiane de cette série de notes est :

	\begin{tblr}{@{}XXXX}
		\textbf{A.} 2&
		\textbf{B.} 3&
		\textbf{C.} 3,5&
		\textbf{D.} 4
	\end{tblr}

\end{document}
